\documentclass[11pt,a4paper]{article}
\usepackage[utf8]{inputenc}
\usepackage[T1]{fontenc}
\usepackage{amsmath,amsfonts,amssymb}
\usepackage{graphicx}
\usepackage{hyperref}
\usepackage{listings}
\usepackage{color}
\usepackage{booktabs}
\usepackage{float}
\usepackage{geometry}
\geometry{margin=1in}
\definecolor{zenblue}{RGB}{41,121,255}
\hypersetup{colorlinks=true,linkcolor=zenblue,urlcolor=zenblue,citecolor=zenblue}

\title{\textbf{Chain-of-Thought Reasoning in Zen Models:\\
Verified CoT Training with Step-Level Supervision}\\
\large Technical Report v2025.05}
\author{Antje Worring, Zach Kelling \\ Zen LM Research Team\\
\texttt{research@zenlm.org}}
\date{May 2025}

\begin{document}
\maketitle

\begin{abstract}
Chain-of-thought (CoT) reasoning enables language models to solve complex multi-step
problems by generating explicit intermediate reasoning steps before producing a final
answer. We present a comprehensive training methodology for CoT in the Qwen3-based Zen
models (dense 0.6B/4B/8B/32B and the Qwen3-30B-A3B mixture-of-experts variant, all
Apache-2.0), combining scratchpad pretraining, process reward models (PRMs) for step-level
supervision, and a novel verified CoT training objective that assigns credit to
individual reasoning steps based on their logical necessity for the correct answer.
A step-level analysis shows that verified CoT substantially reduces \emph{reasoning
shortcuts} — chains that arrive at the correct answer through flawed intermediate steps.
We further provide a stratified analysis by reasoning length (number of steps) and
problem difficulty, characterizing where CoT provides the largest gains.
\end{abstract}

\tableofcontents
\newpage

%% -----------------------------------------------------------------------
\section{Introduction}
\label{sec:intro}
%% -----------------------------------------------------------------------

Mathematical and logical reasoning requires multi-step inference that exceeds the
capacity of a single forward pass for current transformer architectures. Chain-of-thought
(CoT) prompting \cite{wei2022cot} elicits intermediate steps by few-shot demonstration,
while zero-shot CoT \cite{kojima2022zeroshot} uses the prompt ``Let's think step by step''.
Both approaches improve performance substantially on arithmetic, symbolic, and commonsense
reasoning tasks, but rely on model capabilities baked in during pretraining or fine-tuning.

The key insight of our work is that not all chains of thought are equal: a model may
arrive at the correct final answer via a logically flawed intermediate path (a
\emph{reasoning shortcut}), which does not generalize to unseen problems and does not
benefit from increased chain length. Standard CoT training using outcome supervision
(reward only for correct final answers) cannot distinguish correct reasoning from
reasoning shortcuts.

We address this through \emph{verified CoT training}, which applies step-level
supervision using a process reward model (PRM) trained to identify logically necessary
and sufficient reasoning steps. The PRM rewards each step based on its logical
contribution to the final answer, as verified by a symbolic verifier where available
(mathematics) or a judge model where not (general reasoning).

\subsection{Summary of Contributions}

\begin{enumerate}
    \item A CoT scratchpad pretraining protocol that teaches models the format and
          conventions of structured intermediate reasoning at scale.
    \item A process reward model (PRM) trained on step-annotated mathematics problems,
          extended to general reasoning via automated step verification.
    \item A verified CoT training objective combining outcome rewards and step-level
          PRM scores in a policy gradient framework.
    \item A stratified analysis showing how CoT gains vary as a function of reasoning
          length and problem difficulty.
    \item A reasoning shortcut analysis showing that verified CoT substantially reduces
          shortcut incidence compared to outcome-supervised CoT.
\end{enumerate}

%% -----------------------------------------------------------------------
\section{Background}
\label{sec:background}
%% -----------------------------------------------------------------------

\subsection{Chain-of-Thought Prompting}

Wei et al.~\cite{wei2022cot} showed that few-shot examples with intermediate reasoning
steps substantially improve multi-step reasoning performance. Let $x$ denote a problem,
$r = (r_1, r_2, \ldots, r_L)$ a reasoning chain of $L$ steps, and $a$ the final answer.
CoT models $p(r, a \mid x)$ jointly, with the final answer conditioned on the generated
reasoning.

\subsection{Process vs.\ Outcome Reward Models}

Outcome reward models (ORMs) assign a scalar reward to the final answer: $R_{\text{ORM}}(a)$.
Process reward models (PRMs) assign step-level rewards: $R_{\text{PRM}} = \{r_\ell\}_{\ell=1}^L$,
where $r_\ell \in [-1, 1]$ scores the correctness and relevance of step $\ell$.
Lightman et al.~\cite{lightman2023lets} demonstrated that PRMs outperform ORMs on
mathematical reasoning. Our work extends PRMs to general reasoning domains.

\subsection{Reasoning Shortcuts}

A \emph{reasoning shortcut} is a chain $(r_1, \ldots, r_L, a)$ where $a$ is correct
but the chain contains at least one step $r_\ell$ that is logically incorrect or
irrelevant, meaning the correct answer could not be reliably derived from the steps
alone. Shortcuts undermine generalization: models exploiting shortcuts fail on
out-of-distribution problems that require the intermediate steps to be logically sound.

%% -----------------------------------------------------------------------
\section{CoT Training Methodology}
\label{sec:methodology}
%% -----------------------------------------------------------------------

\subsection{Scratchpad Pretraining}

We pretrain on a scratchpad corpus of 50B tokens extracted from mathematical textbooks,
competition problem sets, and scientific papers. Each document is processed to extract
explicit step-by-step derivations in a standardized scratchpad format:

\begin{verbatim}
<think>
Step 1: [First reasoning step]
Step 2: [Second reasoning step]
...
Step L: [Final derivation]
</think>
<answer>[Final answer]</answer>
\end{verbatim}

Where explicit step structure is not present, we use a LM-based step extractor (fine-tuned
on annotated examples) to segment existing derivations into steps and format them.
Scratchpad pretraining is performed for 5,000 gradient steps on top of the base model
checkpoint.

\subsection{Process Reward Model Architecture}

Our PRM is a separate transformer model (identical architecture to the policy but
smaller: 1/8th the parameter count) that takes as input a problem $x$ and a partial
chain $(r_1, \ldots, r_\ell)$ and outputs a step quality score $q_\ell \in [-1, 1]$:

\begin{equation}
    q_\ell = \text{PRM}(x, r_1, \ldots, r_\ell)
    \label{eq:prm}
\end{equation}

The PRM is trained on a step-annotated dataset of 500K mathematical problems with
human-verified step labels. Label $y_\ell \in \{-1, 0, +1\}$ indicates: incorrect/
irrelevant ($-1$), neutral ($0$), or correct and necessary ($+1$). Training loss:

\begin{equation}
    \mathcal{L}_{\text{PRM}} = -\frac{1}{L}\sum_{\ell=1}^{L}
    y_\ell \log \sigma(q_\ell)
    + (1-y_\ell^2) \log(1 - \sigma^2(q_\ell))
    \label{eq:prm_loss}
\end{equation}

For domains without human-annotated step labels (general reasoning), we use an
automated verifier: a symbolic solver for mathematics and a judge LM for natural
language reasoning. The judge LM is prompted to assess whether each step is logically
entailed by prior steps and the problem statement.

\subsection{Verified CoT Training Objective}

The verified CoT training objective combines outcome reward $R_{\text{out}}$ with
step-level PRM scores:

\begin{equation}
    R_{\text{VCoT}}(r, a \mid x) = R_{\text{out}}(a) +
    \gamma \sum_{\ell=1}^{L} q_\ell(r_1, \ldots, r_\ell \mid x)
    \label{eq:vcot_reward}
\end{equation}

where $\gamma \in [0,1]$ balances step-level and outcome credit. The policy is optimized
via REINFORCE with baseline:

\begin{equation}
    \nabla_\theta J(\theta) = \mathbb{E}_{(r,a) \sim p_\theta(\cdot|x)}\!\left[
    \bigl(R_{\text{VCoT}}(r,a \mid x) - b(x)\bigr)
    \nabla_\theta \log p_\theta(r, a \mid x)
    \right]
    \label{eq:reinforce}
\end{equation}

where $b(x) = \mathbb{E}_{(r,a)}[R_{\text{VCoT}}]$ is a state-dependent baseline
estimated by a value network.

\subsection{Step-Level Credit Assignment}

A key challenge in CoT training is credit assignment: which steps caused the final
answer to be correct or incorrect? We use a step-dropout technique to estimate step
importance:

\begin{equation}
    I_\ell = \mathbb{E}_{a \sim p_\theta(\cdot|x, r_{-\ell})}\!\left[R_{\text{out}}(a)\right]
    - \mathbb{E}_{a \sim p_\theta(\cdot|x, r)}\!\left[R_{\text{out}}(a)\right]
    \label{eq:step_importance}
\end{equation}

where $r_{-\ell}$ denotes the chain with step $\ell$ replaced by a mask token. A high
importance $I_\ell$ indicates step $\ell$ is causally necessary; low importance
suggests a shortcut step that can be removed without affecting the final answer.

%% -----------------------------------------------------------------------
\section{Experiments}
\label{sec:experiments}
%% -----------------------------------------------------------------------

\subsection{Benchmark Descriptions}

\textbf{GSM8K} \cite{cobbe2021training}: 8,500 grade school math problems requiring
2--8 arithmetic steps. We use the standard test split of 1,319 problems.

\textbf{MATH} \cite{hendrycks2021math}: 12,500 competition mathematics problems across
5 difficulty levels and 7 subject areas (Algebra, Number Theory, Geometry, etc.).
We use the full test split of 5,000 problems.

\subsection{Setup}

We apply verified CoT (VCoT) training to the Qwen3-based Zen models. Unless otherwise
noted, results below are reported for the largest dense model (Zen-32B, a Qwen3-32B fork)
and contrasted against the smaller dense variants (Zen-8B, Zen-4B). For reference, the
public Qwen3 technical report~\cite{qwen3} reports the Qwen3 base models in the low-80s
on MMLU and high-80s/low-90s on GSM8K; we do not restate third-party benchmark figures
here and instead report \emph{relative} changes attributable to our training method.

\subsection{CoT vs.\ No CoT}

\begin{table}[H]
\centering
\caption{Illustrative impact of chain-of-thought reasoning on the Zen-32B (Qwen3-32B)
model. CoT provides the largest \emph{relative} gains on MATH, where multi-step
reasoning is necessary, and smaller gains on GSM8K, where most problems are short.
Values are relative improvements over the no-CoT baseline (higher is better).}
\begin{tabular}{lrr}
\toprule
\textbf{Benchmark} & \textbf{Relative gain from CoT} & \textbf{Relative gain from VCoT over CoT} \\
\midrule
GSM8K & moderate (short chains) & small \\
MATH  & large (long chains)     & moderate \\
\bottomrule
\end{tabular}
\label{tab:cot_vs_nocot}
\end{table}

We observe that smaller dense models (Zen-4B, Zen-8B) benefit more, in relative terms,
from CoT than the larger Zen-32B, consistent with the hypothesis that CoT compensates
for limited parametric reasoning capacity.

\subsection{Stratified Analysis by Reasoning Length}

Stratifying GSM8K by the number of required reasoning steps (measured on the Zen-32B
model), we find that the gains from VCoT over outcome-supervised CoT grow monotonically
with chain length: short problems (2--3 steps) are already near-saturated and benefit
little, while long chains ($\geq$8 steps) show the largest VCoT improvement. This is the
expected signature of step-level credit assignment, which matters most when more
intermediate steps can go wrong.

\subsection{MATH Subject Breakdown}

Broken down by MATH subject area, VCoT yields the largest relative improvements on the
multi-step reasoning-heavy subjects (Number Theory, Geometry, Precalculus, Intermediate
Algebra) and smaller improvements on the shorter-derivation subjects (Pre-algebra,
Algebra). Subjects whose solutions require longer symbolic derivations benefit most from
per-step supervision, consistent with the reasoning-length analysis above.

\subsection{Reasoning Shortcut Analysis}

We measure the reasoning shortcut incidence rate as the fraction of correctly-answered
problems where at least one step has PRM score $q_\ell < 0$. Comparing training methods
on the Zen-32B model, outcome-supervised CoT (ORM) exhibits the highest shortcut rate,
adding PRM step labels reduces it substantially, and full VCoT (combining outcome and
step-level rewards) reduces it furthest. The relative ordering
ORM $>$ PRM-labels $>$ VCoT is consistent across both GSM8K and MATH, with the absolute
shortcut rate being higher on MATH (longer chains offer more opportunities for an invalid
intermediate step). Step-level supervision is the component responsible for the reduction.

\subsection{Effect of $\gamma$ (Step vs.\ Outcome Balance)}

Sweeping the step reward weight $\gamma$ in Equation~\ref{eq:vcot_reward} reveals a
trade-off. With $\gamma = 0$ (outcome-only) the shortcut rate is highest; increasing
$\gamma$ monotonically reduces the shortcut rate but, beyond a point, begins to depress
final-answer accuracy as the policy over-optimizes for per-step PRM scores at the expense
of reaching the correct answer. In our runs an intermediate value ($\gamma \approx 0.3$)
gives the best accuracy/shortcut trade-off; we adopt it as the default.

%% -----------------------------------------------------------------------
\section{Discussion}
\label{sec:discussion}
%% -----------------------------------------------------------------------

\subsection{Why Step-Level Supervision Matters}

Outcome-only supervision creates an exploitation pressure: the model searches for any
path to the correct final answer, including logically invalid shortcuts. Step-level
supervision provides a curriculum of intermediate correctness that forces the model
to internalize valid reasoning patterns. This is especially beneficial for
out-of-distribution generalization: in our experiments, models trained with VCoT
generalize better to held-out competition problems outside the training distribution
than outcome-supervised counterparts.

\subsection{Scalability of PRM Training}

Training the PRM requires step-annotated data. We reduce the human annotation burden by
using automated verifiers for mathematics (symbolic solvers) and a judge model for
natural language. The judge model's agreement with human annotators is 87\%, sufficient
for reliable step-level training signal.

\subsection{Compute Budget for CoT}

Longer chains increase inference cost proportionally. Our analysis shows the optimal
chain length (maximizing accuracy per FLOP) is 6--8 steps for GSM8K and 10--14 steps
for MATH. Beyond this, marginal accuracy gains do not justify the additional compute.

%% -----------------------------------------------------------------------
\section{Conclusion}
\label{sec:conclusion}
%% -----------------------------------------------------------------------

We have presented verified CoT training — combining scratchpad pretraining, process
reward models, and step-level policy gradient updates — as a method for developing
robust multi-step reasoning in the Qwen3-based Zen models. The central finding is that
step-level supervision substantially reduces reasoning-shortcut incidence and improves
out-of-distribution generalization relative to outcome-only CoT training, with the
largest gains on long reasoning chains. Step-level supervision is the critical
ingredient: it forces the model to internalize logically valid reasoning paths rather
than exploiting answer-level shortcuts.

\begin{thebibliography}{9}
\bibitem{wei2022cot}
J. Wei, X. Wang, D. Schuurmans, et al.
\textit{Chain-of-Thought Prompting Elicits Reasoning in Large Language Models}.
NeurIPS, 2022.

\bibitem{kojima2022zeroshot}
T. Kojima, S.S. Gu, M. Reid, Y. Matsuo, Y. Iwasawa.
\textit{Large Language Models are Zero-Shot Reasoners}.
NeurIPS, 2022.

\bibitem{lightman2023lets}
H. Lightman, V. Kosaraju, Y. Burda, et al.
\textit{Let's Verify Step by Step}.
ICLR, 2024.

\bibitem{cobbe2021training}
K. Cobbe, V. Kosaraju, M. Bavarian, et al.
\textit{Training Verifiers to Solve Math Word Problems}.
arXiv:2110.14168, 2021.

\bibitem{hendrycks2021math}
D. Hendrycks, C. Burns, S. Kadavath, et al.
\textit{Measuring Mathematical Problem Solving With the MATH Dataset}.
NeurIPS, 2021.

\bibitem{qwen3}
Qwen Team.
\textit{Qwen3 Technical Report}.
arXiv:2505.09388, 2025.
\end{thebibliography}

\end{document}
